\section{光的衍射}
\sourceproblem{选择1}{46}{52}
\problemtext 单缝夫琅禾费衍射装置中，单缝沿透镜光轴方向平移，判断焦平面条纹。
\workings 平行光照明、屏位于透镜后焦面时，衍射角分布由缝宽和波长确定，纵向平移只改变传播相位，理想强度图样不变，选C。
\crosscheck 默认透镜有效口径足够且保持焦平面观察。
\sourceproblem{选择2}{46}{52}
\problemtext 固定波长，单缝加宽，中央明纹中心不变，判断各级衍射角。
\workings 暗纹条件 $a\sin\theta=j\lambda$，各级暗纹向中心靠近；整个归一化图样在 $\sin\theta$ 坐标压缩，选A。
\crosscheck 绝对光强通常也改变，不能认为光强恒定。
\sourceproblem{选择3}{46}{52}
\problemtext 波长 $\lambda$ 的光垂直入射，30度方向缝面可分成6个半波带，求缝宽。
\workings $a\sin30^\circ=6\lambda/2=3\lambda$，故 $a=6\lambda$，选B。
\crosscheck 半波带数乘以半波长，而非整个波长。
\sourceproblem{选择4}{46}{52}
\problemtext 图1单缝AB的第一暗纹在Q，求沿出射方向的边缘光程差BC。
\workings 第一暗纹满足 $a\sin\theta=\lambda$，图中BC为两边缘光线的投影光程差，故BC长为 $\lambda$，选B。
\crosscheck 第一暗纹可分成两个相消的半波带。
\sourceproblem{选择5}{46}{52}
\problemtext 波长600 nm、缝宽0.3 mm，中央两侧第一暗纹间距4 mm，求焦距。
\workings 近轴条件下 $w=2f\lambda/a$，故 $f=wa/(2\lambda)=1.0$ m，选C。
\crosscheck 中央明纹全宽包含两侧各一个第一暗纹距离；角度约0.002 rad。
\sourceproblem{选择6}{46}{52}
\problemtext 光栅常数 $a+b$，透光宽度 $b$；第3、6、9级主极大缺级，判断栅距关系。
\workings 主极大 $(a+b)\sin\theta=k\lambda$ 与单缝暗纹 $b\sin\theta=j\lambda$ 重合时缺级。最小正缺级3对应 $(a+b)/b=3$ 的通常选项，即 $a=2b$。
\crosscheck 一般若比值为既约分数 $p/q$，缺级为 $p$ 的倍数；仅凭缺级序列不能排除3/2等非整数比，需结合原选项。
\sourceproblem{选择7}{46}{52}
\problemtext 白光垂直入射光栅，判断同一非零级光谱中离中央最近的颜色。
\workings $d\sin\theta=k\lambda$，同一级下短波长偏角小，紫光最近，选A。
\crosscheck 零级不分色；比较必须在同一级次内。
\sourceproblem{选择8}{46}{52}
\problemtext 选择适于测量单色可见光波长的光栅常数（选项续PDF53）。
\workings 0.01 mm即10微米，比可见光波长大约一个数量级，能产生易测的非零级偏角且保留多个级次，选C。
\crosscheck 0.5及1 mm偏角太小；高清PDF53确认D为 $1.0\times10^{-3}$ mm，即1微米，实际上也可产生可见光一级衍射，但可见级数少且角度较大。选择C依据常规教学测量折中，并非唯一可用栅距。
\sourceproblem{选择9}{47}{53}
\problemtext 光栅常数 $d$、缝宽 $b$、总缝数 $N$，判断主极大条件。
\workings 相邻缝相位差为 $2\pi d\sin\theta/\lambda$，同相叠加要求 $d\sin\theta=k\lambda$，选C。
\crosscheck $N$ 影响主峰宽度，$b$ 决定单缝包络，均不直接替代栅距。
\sourceproblem{选择10}{47}{53}
\problemtext 判断提高衍射受限光学仪器分辨本领的方法。
\workings 圆孔瑞利角分辨极限 $\theta_{\min}\simeq1.22\lambda/D$，增大孔径、减小波长可提高分辨本领，选B。
\crosscheck 实际仪器还可能受像差、噪声等限制。
\sourceproblem{填空1}{47}{53}
\problemtext 400 nm光经单缝，在焦距1.5 m的透镜焦面第三暗纹距中心3.0 mm；另一波长的第二暗纹位于同处，求缝宽及该波长。
\workings 近轴 $a=3\lambda f/x=6.0\times10^{-4}$ m，同角度满足 $3(400)=2\lambda_0$，故 $\lambda_0=600$ nm。
\crosscheck 使用精确角度 $\sin\theta=x/\sqrt{x^2+f^2}$ 时缝宽仅有约百万分之二的相对改动。
\sourceproblem{填空2}{47}{53}
\problemtext 单缝第五级暗纹方向对应多少半波带？缝宽减半后原方向为明纹还是暗纹？
\workings 高清PDF53确认原题为第五级暗纹。$b\sin\theta=5\lambda$，对应10个半波带。缝宽减半后光程差为 $5\lambda/2$，即5个半波带，按教材半波带近似为明纹。
\crosscheck 精确夫琅禾费强度在此处非零，归一化强度为 $4/(25\pi^2)$；次极大真正位置满足 $\tan\beta=\beta$，故此方向并非严格极大点。修正第十八批误读为第二级的问题。
\sourceproblem{填空3}{47}{53}
\problemtext 632.8 nm的氦氖激光单缝衍射，屏距4 m，两侧第二级暗纹间距10 cm，求缝宽。
\workings 每侧距离5 cm，近轴 $a=2\lambda D/0.05=1.01248\times10^{-4}$ m，约0.101 mm。
\crosscheck 无透镜时采用远场近似；以所求缝宽估计菲涅耳数约0.004，近似适用。
\sourceproblem{填空4}{47}{53}
\problemtext 600 nm光正入射宽 $3.0\times10^{-4}$ m 的单缝，中央明纹全宽 $2.0\times10^{-3}$ m，求透镜焦距。
\workings $b\sin\theta_1=\lambda$，$w=2f\tan\theta_1$，故
\[
f=\frac w2\sqrt{(b/\lambda)^2-1}\approx0.500\,\mathrm m.
\]
\crosscheck 高清PDF53确认缝宽指数负4、明纹全宽指数负3；近轴结果为0.5 m，关闭扫描疑点。
\sourceproblem{填空5}{47}{53}
\problemtext 光栅透光宽度与遮光宽度相等，判断可见主极大的可能级次。
\workings $d=2b$，偶数非零级与单缝暗纹重合而缺级，故可能级次为0及正负奇数级，并满足 $|k|\lambda\le d$。
\crosscheck 0级不缺；未给栅距和波长不能确定全部实际可见级次。
\sourceproblem{填空6}{47}{53}
\problemtext 550 nm光正入射栅距 $1.5\times10^{-4}$ cm 的光栅，求最高衍射级。
\workings 栅距1500 nm，$|k|\le1500/550=2.727\ldots$，故最高几何允许级为2。
\crosscheck 若另有偶数缺级结构则实际可见最高级可能更低；题中未给缝宽。
\sourceproblem{填空7}{47}{53}
\problemtext 白光入射光栅，400 nm的第三级与另一波长的第二级重合，求后者。
\workings 同方向 $3(400)=2\lambda_2$，故 $\lambda_2=600$ nm。
\crosscheck 两级必须均未缺级，且栅距允许该角度。

\sourceproblem{计算1}{48}{54}
\problemtext 600 nm光正入射宽0.20 mm单缝，焦距1.0 m，求中央明纹宽度和第二暗纹距焦点距离。
\workings 暗纹满足 $b\sin\theta_j=j\lambda$，焦面坐标 $x_j=f\tan\theta_j$。因此 $\Delta x_0=2f\tan(\arcsin(\lambda/b))\approx6.00\,\mathrm{mm}$，$x_2=f\tan(\arcsin(2\lambda/b))\approx6.00\,\mathrm{mm}$。
\crosscheck 前者是全宽，后者是单侧距离；二者只在小角近似下相等。
\sourceproblem{计算2}{48}{54}
\problemtext 光栅每厘米200缝、缝宽 $2\times10^{-3}$ cm，焦距1 m，入射波长600 nm，求单缝中央包络宽度和其中光栅主明纹数。
\workings $d=5\times10^{-5}$ m，$b=2\times10^{-5}$ m。包络全宽 $2f\tan(\arcsin(\lambda/b))\approx0.0600$ m。主极大位于中央包络内要求 $|k|\lambda/d<\lambda/b$，即 $|k|<2.5$，故 $k=0,\pm1,\pm2$，共5条。
\crosscheck 缝宽指数经高清PDF54确认；第5级等为缺级，但不在中央包络内。
\sourceproblem{计算3}{49}{55}
\problemtext 氦氖激光波长632.8 nm，单缝第一暗纹角5度，求缝宽。
\workings $b=\lambda/\sin5^\circ\approx7.26\times10^{-6}\,\mathrm m$。
\crosscheck 角度较小时用弧度近似，精确结果使用正弦。
\sourceproblem{计算4}{49}{55}
\problemtext 两波长正入射光栅，第一种的三级和第二种的四级主极大均在30度，已知 $\lambda_1=560$ nm，求栅距和第二波长。
\workings $d\sin30^\circ=3\lambda_1=4\lambda_2$，故 $d=3360$ nm，$\lambda_2=420$ nm。
\crosscheck 两级均满足几何允许条件；题设已观察到主极大，故不缺级。
